


In the past few years, several works have discovered the advantage of obtaining segmentation masks in a weakly-supervised manner. In other words, they aim to capture the pixel-level object segmentation mask with image-level labels only. Most of the existing works are initially based on Class Activation Mapping (CAM), which is a technique for identifying activation regions of a trained network. Broadly speaking, recent CAM methods can be categorized into three categories: 1) post-generated CAM, 2) in-training CAM, and 3) refined CAM.


\textbf{Post-generated CAM.} These CAM methods \cite{zhou2016learning,selvaraju2017grad,omeiza2019smooth,wang2020score,ramaswamy2020ablation,jiang2021layercam} train the neural network and generate the activation map separately. They utilize the feature maps from the last convolution layer and generate activation maps for the target object through a weighted sum, where the weight is defined differently. Zhou et al. \cite{zhou2016learning} introduced the very first CAM method, which directly defines the trainable weights of the fully-connected layer as the weight of feature maps. As an extension, Grad-CAM \cite{selvaraju2017grad} was proposed to calculate the gradients between the target class confidence and the feature map and use them as the weights. Different from \cite{selvaraju2017grad}, Grad-CAM can be applied to almost all kinds of convolution neural architectures. Wang et al. \cite{wang2020score} proposed ScoreCAM, which used the logit of the input masked by feature maps to define the weight. Englebert et al. \cite{englebert2022poly} followed a similar idea as ScoreCAM to define the Channel-wise Variation of Confidence as the weight. Besides, they also proposed a recursive step by aggregating the CAM from low resolution to high resolution, which is called Poly-CAM.
Following-up works have discovered the potential of different methods to generate segmentation results from these activation maps in medical images. For example, Nguyen et al. \cite{nguyen2019novel} and Kang et al. \cite{kang2021towards} applied thresholding and refinement on the activation map to obtain the final segmentation result. Alternatively, Jiang et al. \cite{jiang2021layercam} put forward LayerCAM, which applied Grad-CAM to multiple internal layers and add the obtained activation maps to form the final one. 

\textbf{In-training CAM.} Different from post-generated CAM, in-training CAM methods \cite{wang2020self,lee2021lfi,zhang2018adversarial} train the neural network and generate the activation map simultaneously. They reveal that object localization maps can be conveniently obtained by directly selecting from the feature maps of the last convolutional layer through the network itself, which no longer requires an extra step to generate an activation map. 
Zhang et al. \cite{zhang2018adversarial} first proposed to produce class-specific activation maps during the forward pass. In addition, they used two adversary classifiers to find the entire object region, rather than barely focusing on a small region. 
Lee et al. \cite{lee2021lfi} adopted the concept of ScoreCAM to propose the LFI-CAM, which claimed to learn the score more efficiently. Moreover, they leveraged the attention mechanism to improve the classification performance and generated a more consistent attention map during the forward propagation.
Wang et al. \cite{wang2020self} introduced a self-supervised equivariant attention mechanism to keep the consistency of the network for generating activation maps. Similarly, based on Wang et al., Tang et al. \cite{tang2021m} included the concept of multiple-instance learning and neighborhood affinity module to improve the training efficiency and the accuracy of the activation map. 

\textbf{Refined CAM.} Methods \cite{li2022transcam,xie2022c2am} that belong to this category do not directly generate the activation map but provide a clue or a reference map for refinement from initial results.  C$^2$AM \cite{xie2022c2am} is the latest approach that proposed an unsupervised contrastive learning method to generate an activation map of the object. However, they applied the background activation map to the initial CAM obtained from previous works to improve the performance.


As mentioned in Section \ref{sec:intro}, most of these CAM methods generate the object localization map from the low-resolution feature map, which can lead to inferior solutions. LayerCAM partially addresses this problem, but due to the noise in internal activation map extraction and the lack of flexibility in aggregation, often the results are still not ideal. To address these problems, the proposed AME-CAM uses the multiple-exit training strategy to accurately capture the higher-resolution internal activation maps. Moreover, we leverage the attention mechanism to aggregate these activation maps with unsupervised contrastive learning.


